\chapter*{Cómo leer este libro}
\addcontentsline{toc}{chapter}{Cómo leer este libro}

No hace falta leer el libro en orden. Los capítulos~1--4 son la base común: los tipos de neuronas, qué calculan \nt{GLUT} y \nt{GABA} y en qué lenguaje se comunican. A partir de ahí el libro se ramifica. El mapa de la fig.~\ref{fig:roadmap} muestra qué capítulos se apoyan en cuáles: una flecha del capítulo $A$ al capítulo $B$ significa que $B$ usa conceptos y fórmulas de $A$. El mapa recoge solo las dependencias principales; las pequeñas referencias hacia atrás no estorban la lectura.

\begin{figure}[h]
\centering
\begin{tikzpicture}[>=Stealth,scale=0.95,
  ch/.style={draw,thick,rounded corners=3pt,align=center,font=\scriptsize,text width=1.85cm,minimum height=0.62cm,inner sep=2pt},
  base/.style={ch,fill=blue!8}, bus/.style={ch,fill=orange!14}, sum/.style={ch,fill=gray!18},
  io/.style={ch,fill=green!10}, sort/.style={ch,fill=cyan!8}, lab/.style={ch,fill=violet!10},
  dep/.style={->,thick,gray!60!black}]
% foundation
\node[base] (c1) at (0,0) {1 Tipos de neuronas};
\node[base] (c2) at (2.3,0) {2 \nt{GLUT}};
\node[base] (c3) at (4.6,0) {3 \nt{GLUT} y \nt{GABA}};
\node[base] (c4) at (6.9,0) {4 Código Morse};
% broadcast channels and the synthesis
\node[bus] (c5) at (4.6,-1.5) {5 \nt{SERO}};
\node[bus] (c6) at (7.3,-1.5) {6 \nt{DOPA}};
\node[bus] (c7) at (9.2,-0.9) {7 Teorías de la dopamina};
\node[bus] (c8) at (9.2,-2.1) {8 \nt{NORA}};
\node[bus] (c9) at (11.5,-2.1) {9 \nt{ACET}};
\node[sum] (c10) at (13.8,-0.9) {10 La máquina completa};
% inputs and outputs
\node[io] (c12) at (2.3,-4.1) {12 Recono-\\cimiento};
\node[io] (c11) at (4.6,-4.1) {11 Entradas};
\node[io] (c13) at (6.9,-4.1) {13 Músculos};
\node[io] (c14) at (9.2,-3.05) {14 Dolor};
\node[io] (c15) at (9.2,-4.1) {15 Aprender\\movimientos};
\node[io] (c16) at (9.2,-5.15) {16 Química\\del cuerpo};
% lab, sorting, sleep
\node[lab] (c17) at (11.5,-4.1) {17\\Depuración};
\node[lab] (c21) at (13.8,-4.1) {21 Máquinas vivas};
\node[lab] (c22) at (13.8,-5.15) {22 DishBrain};
\node[sort] (c18) at (11.5,-5.15) {18 Neuronas-\\instrumento};
\node[sort] (c19) at (13.8,-6.2) {19 Número\\de dendritas};
\node[bus] (c20) at (11.5,-6.2) {20 Sueño};
% dependencies
\draw[dep] (c1) -- (c2); \draw[dep] (c2) -- (c3); \draw[dep] (c3) -- (c4);
\draw[dep] (c1) -- (c5); \draw[dep] (c4) -- (c6); \draw[dep] (c5) -- (c6);
\draw[dep] (c6) -- (c7); \draw[dep] (c6) -- (c8); \draw[dep] (c8) -- (c9);
\draw[dep] (c7) -- (c10); \draw[dep] (c9) -- (c10);
\draw[dep] (c4) -- (c11); \draw[dep] (c11) -- (c12); \draw[dep] (c11) -- (c13); \draw[dep] (c5) -- (c13);
\draw[dep] (c13) -- (c14); \draw[dep] (c13) -- (c15); \draw[dep] (c13) -- (c16);
\draw[dep] (c15) -- (c17); \draw[dep] (c9) -- (c17); \draw[dep] (c17) -- (c21); \draw[dep] (c21) -- (c22);
\draw[dep] (c16) -- (c18); \draw[dep] (c18) -- (c19); \draw[dep] (c16) -- (c20);
% legend
\begin{scope}[yshift=-7.25cm]
\foreach \s/\t [count=\i] in {base/base, bus/canales y modos, sum/síntesis, io/entradas y salidas, sort/clasificaciones de neuronas, lab/laboratorio}{
  \pgfmathtruncatemacro{\col}{mod(\i-1,3)}\pgfmathtruncatemacro{\row}{(\i-1)/3}
  \node[\s,text width=0.3cm,minimum height=0.32cm] at ({\col*4.8+0.2},{-\row*0.55}) {};
  \node[anchor=west,font=\scriptsize] at ({\col*4.8+0.55},{-\row*0.55}) {\t};}
\end{scope}
\end{tikzpicture}
\caption{Mapa de capítulos. La flecha va de un capítulo al que se apoya en él. Las demás referencias entre capítulos son indicaciones, no dependencias.}
\label{fig:roadmap}
\end{figure}

\section*{Recorridos por el libro}

\begin{center}
\footnotesize
\setlength{\tabcolsep}{4pt}
\renewcommand{\arraystretch}{1.2}
\begin{tabular}{>{\raggedright}p{3.0cm}>{\raggedright}p{3.9cm}>{\raggedright\arraybackslash}p{6.4cm}}
\toprule
Recorrido & Para quién & Capítulos en orden \\
\midrule
curso de un trimestre & docente, 10 semanas & semana 1: caps.~1--2; 2: cap.~3; 3: cap.~4; 4: caps.~5--6; 5: caps.~7--8; 6: cap.~9; 7: cap.~10; 8: caps.~11--12; 9: caps.~13 y~15; 10: cap.~17 \\
IA y aprendizaje automático & quien conoce las redes neuronales y quiere entender cómo aprende el cerebro & 1--4, 5 (lectura rápida), 6, 7, 8, 9, 11 (lectura rápida), 12, 13 (lectura rápida), 15 \\
electrónica y hardware & diseñador de circuitos, desarrollador de chips neuromórficos & 1--4, 5, 11, 13, 16, 18, 19, 17 (caps.~9 y 15 en lectura rápida) \\
neurointerfaces y computadoras vivas & quien construye interfaces cerebro-computadora y chips con neuronas vivas & 1, 2, 4, 11, 13, 17 (caps.~9 y 15 en lectura rápida), 21, 22 (cap.~12 en lectura rápida) \\
panorama rápido & ingeniero con un fin de semana libre & 1, 2, 3, 6, 10; las referencias del cap.~6 a los caps.~4--5 se entienden por el contexto \\
\bottomrule
\end{tabular}
\end{center}

“Lectura rápida” significa que basta con leer el comienzo del capítulo, sus proposiciones en los recuadros amarillos y la tabla final.

Al final de cada capítulo hay ejercicios: estimaciones, deducciones, simulación y diseño; las respuestas breves están en el apéndice~\ref{sec:answers}. Toda la notación está en el apéndice~\ref{sec:notation}, y los experimentos en que se apoyan las proposiciones principales de cada capítulo, en el apéndice~\ref{sec:supplement}. Los números del libro se obtienen con pequeños programas de la carpeta \texttt{models/}; se pueden ejecutar y modificar.
